% !TEX program = xelatex
% GENERATED by tools/build.py from data/records.csv and templates/publication_list.tex.
% Every entry here also appears, with the same key, in references.bib.
\documentclass[11pt, a4paper]{cvClass}
\geometry{left=2cm, top=1.4cm, right=2cm, bottom=1.6cm, footskip=.7cm}
% ---------------------------------------------------------------------------
% Shared layout macros for cv_long.tex, cv_short.tex and publication_list.tex
% (input by each template; lives in templates/, copied next to the .tex)
% ---------------------------------------------------------------------------
\colorlet{awesome}{awesome-red}          % the CV orange-red (#DC3522)
\definecolor{accent2}{HTML}{E3B23C}      % muted yellow, used sparingly
\setlength{\parskip}{0pt}
\urlstyle{same}

% A hanging list: year (or label) on the left, text on the right
\newlength{\cvlabelwidth}\setlength{\cvlabelwidth}{2.45cm}
\newlength{\cvitemsep}\setlength{\cvitemsep}{1.4mm}
\newenvironment{cvlist}{%
  \vspace{1.2mm}%
  \begin{list}{}{%
    \setlength{\leftmargin}{\cvlabelwidth}\setlength{\labelwidth}{\dimexpr\cvlabelwidth-3mm\relax}%
    \setlength{\labelsep}{3mm}\setlength{\itemsep}{\cvitemsep}%
    \setlength{\parsep}{0pt}\setlength{\topsep}{0pt}\setlength{\partopsep}{0pt}}%
  \fontsize{9pt}{11.5pt}\bodyfontlight\color{text}\raggedright}%
  {\end{list}\vspace{1mm}}
\newcommand{\cvitem}[2]{\item[{\fontsize{8.5pt}{10pt}\bodyfont\color{graytext}#1\hfill}]#2}
\newcommand{\skilllabel}[1]{{\bodyfont\color{darktext}#1}}
\renewcommand{\textbf}[1]{{\bodyfont\bfseries\color{darktext}#1}}


\name{Tobias}{Stål}
\position{Publication list}
\address{37 entries · updated 8 October 2026}
\email{tobias.staal@utas.edu.au}
\homepage{tobbetripitaka.github.io/Tobias\_Staal}
\extrainfo{\href{https://orcid.org/0000-0002-4323-6748}{ORCID 0000-0002-4323-6748}}
\makecvfooter{}{Tobias Stål · Publication list}{\thepage}

% One reference: key in the margin, full citation, links, abstract
\newenvironment{refentry}[1]{%
  \par\vspace{2.2mm}\noindent
  \begin{minipage}[t]{\linewidth}\raggedright\fontsize{9.5pt}{12pt}\bodyfontlight\color{text}%
  {\fontsize{7pt}{8pt}\bodyfont\color{awesome}#1}\\[0.4mm]}%
  {\end{minipage}\par}
\newcommand{\refauthors}[1]{{\bodyfont\color{darktext}#1}}
\newcommand{\reftitle}[1]{{\bodyfont\bfseries\color{darktext}#1}}
\newcommand{\refabstract}[1]{\par\vspace{0.8mm}{\fontsize{8pt}{10pt}\bodyfontlight\color{graytext}#1\par}}

\begin{document}
\makecvheader[C]
\par\vspace{4mm}
{\fontsize{9pt}{11pt}\bodyfontlight\color{graytext}
Complete list with full author names and, where available, abstracts. Grey keys refer to entries in
\href{https://tobbetripitaka.github.io/Tobias\_Staal/publications/references.bib}{references.bib}, where every item can be copied in BibTeX format.
Within each category, newest first.\par}

\cvsection{Peer-reviewed publications}

\begin{refentry}{halpin2026ediacaran}
\refauthors{Jacqueline A. Halpin, Nathan R. Daczko, Jacob A. Mulder, Alessandro Maritati, Tobias Stål} (2026). \reftitle{Ediacaran–Cambrian High‐Strain Melt‐Present Deformation of the Archean Charcot Province, East Antarctica}. \textit{Tectonics}, 45(8), e2026TC009479. American Geophysical Union (AGU).\\
\href{https://doi.org/10.1029/2026TC009479}{doi:10.1029/2026TC009479}
\refabstract{The assembly of Gondwana involved several Pan‐African orogens, although their extension into ice‐covered Antarctica remains contentious. Particularly problematic is our understanding of the Kuunga orogenic system in East Antarctica, including its eroded architecture and temporal evolution. The remote Charcot Province, at the easternmost extent of the Antarctic Kuunga Orogen, presents a key opportunity to test recent tectonic models. Here we investigate the geological history of Alligator Island, a rarely visited exposure of the Charcot Province. Field relationships, petrography and zircon U–Pb–Hf geochronology indicate that the Alligator Island migmatitic gneisses originated from a Mesoarchean (c. 2.97–2.95 Ga) volcano‐sedimentary package. Metamorphism during the Neoarchean (c. 2.76 Ga) was characterized by channeled melt migration. An apparent tectonic quiescence followed until the Ediacaran–Cambrian (580–540 Ma) when high‐strain melt‐present deformation led to pervasive migmatitic textures and tight folding. Zircon Hf isotopic signatures from anatectic 580–540 Ma grains indicate that melt was sourced from crust older than 3.0 Ga. Melt metasomatism affected whole rock chemistry and feldspar Pb–Pb compositions highlighting complex interactions within the crust. Correlations within the Charcot Province suggest that Alligator Island's Mesoarchean volcano‐sedimentary package was deposited on a basement complex possibly represented by Paleo‐Mesoarchean gneisses exposed 40–60 km to the east at Cape Charcot and Davis Peninsula. Overall our findings are consistent with a recent model in which Charcot Province rocks formed through gravitational spreading within the Kuunga orogenic system during Gondwana assembly, and enhance our understanding of the tectonic processes shaping East Antarctica.}
\end{refentry}

\begin{refentry}{manassero2026bed}
\refauthors{Maria Constanza Manassero, Kate Selway, Tobias Stål, Matthias Scheiter, Mareen Loesing, Felicity McCormack, Jacqueline A. Halpin, Bernd Kulessa, Anya M. Reading} (2026). \reftitle{Bed Topography and Subglacial Conditions of Denman Glacier, East Antarctica: Insights From Magnetotelluric Data and Interdisciplinary Studies}. \textit{Journal of Geophysical Research: Earth Surface}, 131(7), e2026JF009277. American Geophysical Union (AGU).\\
\href{https://doi.org/10.1029/2026JF009277}{doi:10.1029/2026JF009277}
\refabstract{Understanding the basal conditions of East Antarctica's outlet glaciers is critical for accurately modeling future ice‐sheet evolution. Denman Glacier, one of the continent's fastest‐flowing glaciers, has been identified in bed‐topography models as occupying the deepest continental trough on Earth, potentially acting as a major outlet of the East Antarctic Ice Sheet. Here we present a land‐based magnetotelluric transect across Denman Glacier, acquired approximately 50 km upstream from its grounding line. These measurements provide the first in situ, ground‐based geophysical constraints on the glacier's subglacial structure, enabling evaluation of existing airborne‐derived models and reducing uncertainty and non‐uniqueness in inferred bed topography. The resistivity model images a laterally continuous, highly resistive unit (–), interpreted as glacial ice. This ice overlies lower‐resistivity regions (–) interpreted to be unconsolidated sediments and rock. Along the glacier flanks, the boundary between ice and underlying material aligns closely with radar‐derived trough depths. In contrast, beneath the glacier center, the model indicates a trough depth of m below sea level, more than 1000 m shallower than previous estimates, indicating that the extreme depths predicted by earlier models are not observed. Basement resistivity is consistent with magnetic source depths and regional geological constraints. A low‐resistivity feature () near the center of the glacier is interpreted as a sedimentary topographic high. Low‐resistivity anomalies beneath the glacier are attributed to saline porewater within till or sedimentary rocks. The inferred bed geometry and basal conditions provide improved constraints for modeling Denman Glacier's stability and future evolution.}
\end{refentry}

\begin{refentry}{kelly2026sediment}
\refauthors{I. D. Kelly, A. M. Reading, T. Stål, B. Kulessa, A. García‐Jerez, J. Piña‐Flores, E. Paolucci, A. Tanzini, R. J. Turner, J. C. Magyar, A. P. Bassom} (2026). \reftitle{Determining the Character of Subglacial Sediments in the Ice‐Bedrock Interface Zone of Antarctica Using Horizontal‐to‐Vertical Spectral Ratios (HVSRs) of Seismic Ambient Noise}. \textit{Journal of Geophysical Research: Solid Earth}, 131(6), e2025JB033029. American Geophysical Union (AGU).\\
\href{https://doi.org/10.1029/2025JB033029}{doi:10.1029/2025JB033029}
\refabstract{Interactions between the Antarctic Ice Sheet and the underlying solid Earth occur within the ice‐bedrock interface zone (IBIZ), containing structures of sediments and rocks that strongly influence ice sheet dynamics. Existing insights into the Antarctic IBIZ come primarily from interpretations of airborne geophysical data and active seismic measurements, which require significant logistics. Passive seismic methods provide a lower cost, lower resolution alternative for subglacial mapping, with the potential for continuous remote monitoring. Here, we utilize horizontal‐to‐vertical spectral ratios (HVSRs) of seismic ambient noise to infer the presence of subglacial low‐velocity zones generated by unlithified sediments or porous, water saturated sedimentary rocks, using existing data from over 80 broadband stations across Antarctica. From 1‐D forward modeling, we identify a decrease in synthetic HVSR peak frequencies dependent on the thickness and S‐wave velocity () of the low‐velocity zone unrelated to seismic wavefield conditions. This sensitivity in HVSR peak frequencies is validated by the agreement between synthetic and observed HVSRs at Antarctic stations with independent prior constraints on subglacial structure. Our observed HVSR analysis additionally reveals a widespread seasonality in observed HVSR amplitudes between 0.2 and 1 Hz that correlates with sea ice evolution and the modulation of secondary microseism energies in the Southern Ocean. From the decrease in observed HVSR peak frequencies, we infer the likelihood of subglacial low‐velocity zones across Antarctica along major passive seismic transects, mapping to key subglacial sedimentary basins. We provide practical recommendations for future HVSR applications to infer subglacial low‐velocity zones reliably.}
\end{refentry}

\begin{refentry}{mccormack2026synthetic}
\refauthors{Felicity S. McCormack, Tobias Stål, Niya Shao, Emma MacKie, Ana Fabela Hinojosa, Mareen Lösing, Jason Roberts, Shivani Ehrenfeucht, Christine Dow} (2026). \reftitle{Synthetic bed topographies for Antarctica and their utility in ice sheet modelling}. \textit{Philosophical Transactions of the Royal Society A: Mathematical, Physical and Engineering Sciences}, 384(2319), 20240537. The Royal Society.\\
\href{https://doi.org/10.1098/rsta.2024.0537}{doi:10.1098/rsta.2024.0537}
\refabstract{Bed topography is a key control on the evolution of the Antarctic Ice Sheet, influencing ice flow, grounding line retreat and the rate and timing of ice mass loss. To assess the sensitivity of ice sheet evolution to bed variability in ice sheet models, synthetic gridded bed topography datasets are often used. Here, we review methods commonly used to generate synthetic beds, their associated uncertainties and the influence of the approach on the characteristics of the resulting bed. Using the Aurora Subglacial Basin in East Antarctica as a case study, we evaluate the impact of five synthetic bed generation methods on projected ice mass loss under a high emission scenario. Sea-level rise estimates vary by up to 11\% (SSP5-8.5 forcing scenario) and 32\% (RCP2.6) at 2300 CE when basal friction coefficients from the friction law are optimized for each bed, and by up to 23\% (SSP5-8.5) and 51\% (RCP2.6) at 2300 CE when using non-optimized coefficients. Our results highlight the importance of relatively small bed variations on the timing and extent of grounding line retreat and the need for process-informed representation of the basal friction in decadal- to centennial-scale sea-level projections. This article is part of the Theo Murphy meeting issue ‘Next generation ice-sheet bed measurements’.}
\end{refentry}

\begin{refentry}{arora2026geothermal}
\refauthors{Devsamridhi Arora, Naresh Chandra Pant, Tobias Stål, Aditya Naik, Mayuri Pandey, Rashmi Gupta} (2026). \reftitle{Geothermal heat flow in the Princess Elizabeth Land sector, East Antarctica}. \textit{Earth-Science Reviews}, 277, 105475. Elsevier BV.\\
\href{https://doi.org/10.1016/j.earscirev.2026.105475}{doi:10.1016/j.earscirev.2026.105475}
\end{refentry}

\begin{refentry}{alaghbary2026uncertainty}
\refauthors{Magued Al-Aghbary, Mohamed Sobh, Tobias Stål, Mohamed Osman Awaleh, Mohamed Jalludin, Christian Gerhards} (2026). \reftitle{Uncertainty quantification using clustering-based quantile regression forests: with an application to improving geothermal heat flow predictions}. \textit{Geophysical Journal International}, 245(2), ggag113. Oxford University Press (OUP).\\
\href{https://doi.org/10.1093/gji/ggag113}{doi:10.1093/gji/ggag113}
\refabstract{Machine learning models offer powerful predictive capabilities for geoscientific applications but remain limited by their ‘black-box’ nature and lack of rigorous uncertainty quantification. We developed a comprehensive, generalizable uncertainty quantification framework that decomposes predictive uncertainty into aleatoric and epistemic components using Quantile Regression Forests. Additionally, we applied unsupervised k-means clustering to isolate homogeneous data regimes, thereby reducing aleatoric uncertainty across spatially heterogeneous geoscientific data sets. To facilitate interpretation and quality assessment, we introduced five spatial diagnostic tools: bandwidth, variance, robustness, confidence and explainability maps that characterize prediction reliability and identify dominant uncertainty sources. To demonstrate the framework’s applicability, we tested it on three synthetic data sets varying in size and a real-world geothermal heat flow application with 14 geophysical observables across continental Africa. Results show that clustering substantially reduces aleatoric uncertainty while maintaining stable epistemic uncertainty. Clustering also improves predictive accuracy and sharpens prediction intervals, with gains most pronounced in homogeneous regions. Applied to the African geothermal heat flow, the framework reveals region-specific geological controls (lithospheric architecture dominates stable cratons, while tectonic proximity governs active rift zones) and guides targeted data collection by distinguishing high-epistemic regions requiring additional sampling from high-aleatoric zones needing improved observables. While theoretically applicable to other geographic regions and geophysical data sets, the framework’s performance in different geological settings requires validation. This interpretable, uncertainty-aware approach enhances trustworthiness of predictions in spatially heterogeneous, data-sparse geoscientific problems.}
\end{refentry}

\begin{refentry}{losing2026community}
\refauthors{Mareen Lösing, William Colgan, Tobias Stål, Jörg Ebbing, Anne G. Busck, Tong Zhang, Hélène L. Seroussi, Felicity McCormack, Dominik Fahrner, Leigh Stearns, Synne H. Svendsen, Anya Reading} (2026). \reftitle{Community heat flow recommendations: suitable basal boundary conditions for Greenland and Antarctica in ISMIP7}. \textit{GEUS Bulletin}, 62. Geological Survey of Denmark and Greenland.\\
\href{https://doi.org/10.34194/r0w9rf81}{doi:10.34194/r0w9rf81} \quad \url{https://geusbulletin.org/index.php/geusb/article/view/8411}
\refabstract{Geothermal heat flow (GHF) influences ice sheet thermal conditions, affecting ice flow by sliding and deformation. However, GHF distribution under polar ice sheets remains poorly constrained, with few direct borehole-derived estimates and large discrepancies between glaciological and geophysical models caused by methodological differences and data limitations. As a result, many ice sheet models rely on uniform GHF estimates, ensemble averages or outdated fields that oversimplify reality. The choice of GHF product can lead to significantly different thermal conditions simulated at the ice-bed interface, which affects the projected evolution of ice sheets under climate warming. Therefore, we conducted an expert elicitation survey to identify the most suitable GHF fields for use as basal boundary conditions in ice sheet modelling, particularly for the Ice Sheet Modelling Intercomparison Project for CMIP7 (ISMIP7). GHF fields generally fall into three categories: (1) outdated due to improved data availability, (2) overly simplified parameterisations and (3) current and preferred. For GHF fields that rank highly in the survey, we discuss uncertainty and data dependency and guide their use in different applications. Finally, we recommend two Antarctic and one Greenlandic GHF fields for ISMIP7.}
\end{refentry}

\begin{refentry}{losing2025linking}
\refauthors{Mareen Lösing, Alan Aitken, Jörg Ebbing, Jacqueline A. Halpin, Lu Li, Max Moorkamp, Anya Reading, Tobias Stål} (2025). \reftitle{Linking Tectonics and Crustal Thermal Properties in Southwestern Australia and East Antarctica Through Coupled Gravity and Magnetic Analysis}. \textit{Journal of Geophysical Research: Solid Earth}, 130(12), e2024JB030770. American Geophysical Union (AGU).\\
\href{https://doi.org/10.1029/2024JB030770}{doi:10.1029/2024JB030770}
\refabstract{The shared tectonic history of southwestern Australia and East Antarctica facilitates the exchange of geological insights between the regions. In this study, we present coupled susceptibility and density models obtained through the joint inversion of magnetic and gravity data. By assuming a common geological source for both signals, our coupling method minimizes misfits and variation in information, thereby enhancing a correlation between susceptibility and density. The resulting anomalies demonstrate structural continuity between the continents, aligning closely with major shear zones and seismic reflectors. Combining these results with machine learning, geochemical, and petrophysical databases, we predict a high‐resolution (10 km) heat production map for East Antarctica. Utilizing a Markov Chain Monte Carlo (MCMC) algorithm, we further develop a geothermal heat flow map with greater spatial variability than previous studies, yielding an average of mW/ in East Antarctica and mW/ in southwestern Australia. Our results provide a crucial high‐resolution boundary condition for ice sheet simulations, enabling more realistic estimates of basal meltwater production and ice temperatures.}
\end{refentry}

\begin{refentry}{abram2025emerging}
\refauthors{Nerilie J. Abram, Ariaan Purich, Matthew H. England, Felicity S. McCormack, Jan M. Strugnell, Dana M. Bergstrom, Tessa R. Vance, Tobias Stål, Barbara Wienecke, Petra Heil, Edward W. Doddridge, Jean-Baptiste Sallée, Thomas J. Williams, Anya M. Reading, Andrew Mackintosh, Ronja Reese, Ricarda Winkelmann, Ann Kristin Klose, Philip W. Boyd, Steven L. Chown, Sharon A. Robinson} (2025). \reftitle{Emerging evidence of abrupt changes in the Antarctic environment}. \textit{Nature}, 644(8077), 621–633. Springer Science and Business Media LLC.\\
\href{https://doi.org/10.1038/s41586-025-09349-5}{doi:10.1038/s41586-025-09349-5}
\end{refentry}

\begin{refentry}{staal2024geology}
\refauthors{Tobias Stål, Jacqueline A. Halpin, John W. Goodge, Anya M. Reading} (2024). \reftitle{Geology Matters for Antarctic Geothermal Heat}. \textit{Geophysical Research Letters}, 51(13), e2024GL110098. American Geophysical Union (AGU).\\
\href{https://doi.org/10.1029/2024GL110098}{doi:10.1029/2024GL110098}
\refabstract{Geothermal heat plays a vital role in Antarctic ice sheet stability. The continental geothermal heat flow distribution depends on lithospheric composition and ongoing tectonism. Heat‐producing elements are unevenly enriched in the crust over deep time by various geological processes. The contribution of crustal heat production to geothermal heat flow is widely recognized; however, in Antarctica, crustal geology is largely hidden, and its complexity has frequently been excluded in thermal studies due to limited observations and oversimplified assumptions. Li and Aitken (2024), https://doi.org/10.1029/2023GL106201 take a significant step forward, focusing on Antarctic crustal radiogenic heat. Utilizing gravity inversion and rock composition data, they show that the crustal heterogeneity introduces considerable variability to heat flow. However, modeling crustal heat production proves challenging because it lacks distinct associations with geophysical observables and has a narrow spatial association. Robust quantification of geothermal heat production and heat flow must incorporate explicit aspects of geology.}
\end{refentry}

\begin{refentry}{mccormack2022fine}
\refauthors{Felicity S. McCormack, Jason L. Roberts, Christine F. Dow, Tobias Stål, Jacqueline A. Halpin, Anya M. Reading, Martin J. Siegert} (2022). \reftitle{Fine‐Scale Geothermal Heat Flow in Antarctica Can Increase Simulated Subglacial Melt Estimates}. \textit{Geophysical Research Letters}, 49(15), e2022GL098539. American Geophysical Union (AGU).\\
\href{https://doi.org/10.1029/2022GL098539}{doi:10.1029/2022GL098539}
\refabstract{Antarctic geothermal heat flow (GHF) affects the thermal regime of ice sheets and simulations of ice and subglacial meltwater discharge to the ocean, but remains poorly constrained. We use an ice sheet model to investigate the impact of GHF anomalies on subglacial meltwater production in the Aurora Subglacial Basin, East Antarctica. We find that spatially‐variable GHF fields produce more meltwater than a constant GHF with the same background mean, and meltwater production increases as the resolution of GHF anomalies increases. Our results suggest that model simulations of this region systematically underestimate meltwater production using current GHF models. We determine the minimum basal heating required to bring the basal ice temperature to the pressure melting point, which should be taken together with the scale‐length of likely local variability in targeting in‐situ GHF field campaigns.}
\end{refentry}

\begin{refentry}{staal2022properties}
\refauthors{Tobias Stål, Anya M. Reading, Sven Fuchs, Jacqueline A. Halpin, Mareen Lösing, Ross J. Turner} (2022). \reftitle{Properties and biases of the global heat flow compilation}. \textit{Frontiers in Earth Science}, 10, 963525. Frontiers Media SA.\\
\href{https://doi.org/10.3389/feart.2022.963525}{doi:10.3389/feart.2022.963525}
\refabstract{Geothermal heat flow is inferred from the gradient of temperature values in boreholes or short-penetration probe measurements. Such measurements are expensive and logistically challenging in remote locations and, therefore, often targeted to regions of economic interest. As a result, measurements are not distributed evenly. Some tectonic, geologic and even topographic settings are overrepresented in global heat flow compilations; other settings are underrepresented or completely missing. These limitations in representation have implications for empirical heat flow models that use catalogue data to assign heat flow by the similarity of observables. In this contribution, we analyse the sampling bias in the Global Heat Flow database of the International Heat Flow Commission; the most recent and extensive heat flow catalogue, and discuss the implications for accurate prediction and global appraisals. We also suggest correction weights to reduce the bias when the catalogue is used for empirical modelling. From comparison with auxiliary variables, we find that each of the following settings is highly overrepresented for heat flow measurements; continental crust, sedimentary rocks, volcanic rocks, and Phanerozoic regions with hydrocarbon exploration. Oceanic crust, cratons, and metamorphic rocks are underrepresented. The findings also suggest a general tendency to measure heat flow in areas where the values are elevated; however, this conclusion depends on which auxiliary variable is under consideration to determine the settings. We anticipate that using our correction weights to balance disproportional representation will improve empirical heat flow models for remote regions and assist in the ongoing assessment of the Global Heat Flow database.}
\end{refentry}

\begin{refentry}{reading2022antarctic}
\refauthors{Anya M. Reading, Tobias Stål, Jacqueline A. Halpin, Mareen Lösing, Jörg Ebbing, Weisen Shen, Felicity S. McCormack, Christine S. Siddoway, Derrick Hasterok} (2022). \reftitle{Antarctic geothermal heat flow and its implications for tectonics and ice sheets}. \textit{Nature Reviews Earth \& Environment}, 3(12), 814–831. Springer Science and Business Media LLC.\\
\href{https://doi.org/10.1038/s43017-022-00348-y}{doi:10.1038/s43017-022-00348-y}
\end{refentry}

\begin{refentry}{morse2022exploratory}
\refauthors{Peter E. Morse, Anya M. Reading, Tobias Stal} (2022). \reftitle{Exploratory Volumetric Deep Earth Visualization by 2.5D Interactive Compositing}. \textit{IEEE Transactions on Visualization and Computer Graphics}, 28(7), 2641–2653. Institute of Electrical and Electronics Engineers (IEEE).\\
\href{https://doi.org/10.1109/TVCG.2020.3037226}{doi:10.1109/TVCG.2020.3037226}
\end{refentry}

\begin{refentry}{staal2021antarctic}
\refauthors{Tobias Stål, Anya M. Reading, Jacqueline A. Halpin, Joanne M. Whittaker} (2021). \reftitle{Antarctic Geothermal Heat Flow Model: Aq1}. \textit{Geochemistry, Geophysics, Geosystems}, 22(2), e2020GC009428. American Geophysical Union (AGU).\\
\href{https://doi.org/10.1029/2020GC009428}{doi:10.1029/2020GC009428}
\refabstract{We present a refined map of geothermal heat flow for Antarctica, Aq1, based on multiple observables. The map is generated using a similarity detection approach by attributing observables from geophysics and geology to a large number of high‐quality heat flow values ( N = 5,792) from other continents. Observables from global, continental, and regional datasets for Antarctica are used with a weighting function that allows the degree of similarity to increase with proximity and how similar the observables are. The similarity detection parameters are optimized through cross correlation. For each grid cell in Antarctica, a weighted average heat flow value and uncertainty metrics are calculated. The Aq1 model provides higher spatial resolution in comparison to previous results. High heat flow is shown in the Thwaites Glacier region, with local values over 150 mW m −2 . We also map elevated values over 80 mW m −2 in Palmer Land, Marie Byrd Land, Victoria Land and Queen Mary Land. Very low heat flow is shown in the interior of Wilkes Land and Coats Land, with values under 40 mW m −2 . We anticipate that the new geothermal heat flow map, Aq1, and its uncertainty bounds will find extended use in providing boundary conditions for ice sheet modeling and understanding the interactions between the cryosphere and solid Earth. The computational framework and open architecture allow for the model to be reproduced, adapted and updated with additional data, or model subsets to be output at higher resolution for regional studies.}
\end{refentry}

\begin{refentry}{sanchez2021petrochron}
\refauthors{G. Sanchez, J. A. Halpin, M. Gard, D. Hasterok, T. Stål, T. Raimondo, S. Peters, A. Burton‐Johnson} (2021). \reftitle{PetroChron Antarctica: A Geological Database for Interdisciplinary Use}. \textit{Geochemistry, Geophysics, Geosystems}, 22(12), e2021GC010154. American Geophysical Union (AGU).\\
\href{https://doi.org/10.1029/2021GC010154}{doi:10.1029/2021GC010154}
\refabstract{We present PetroChron Antarctica, a new relational database including petrological, geochemical and geochronological data sets along with computed rock properties from geological samples across Antarctica. The database contains whole‐rock geochemistry with major/trace element and isotope analyses, geochronology from multiple isotopic systems and minerals for given samples, as well as an internally consistent rock classification based on chemical analysis and derived rock properties (i.e., chemical indices, density, p ‐velocity, and heat production). A broad range of meta‐information such as geographic location, petrology, mineralogy, age statistics and significance are also included and can be used to filter and assess the quality of the data. Currently, the database contains 11,559 entries representing 10,056 unique samples with varying amounts of geochemical and geochronological data. The distribution of rock types is dominated by mafic (36\%) and felsic (33\%) compositions, followed by intermediate (22\%) and ultramafic (9\%) compositions. Maps of age distribution and isotopic composition highlight major episodes of tectonic and thermal activity that define well known crustal heterogeneities across the continent, with the oldest rocks preserved in East Antarctica and more juvenile lithosphere characterizing West Antarctica. PetroChron Antarctica allows spatial and temporal variations in geology to be explored at the continental scale and integrated with other Earth‐cryosphere‐biosphere‐ocean data sets. As such, it provides a powerful resource ready for diverse applications including plate tectonic reconstructions, geological/geophysical maps, geothermal heat flow models, lithospheric and glacial isostasy, geomorphology, ice sheet reconstructions, biodiversity evolution, and oceanography.}
\end{refentry}

\begin{refentry}{staal2020antarctic}
\refauthors{Tobias Stål, Anya M. Reading, Jacqueline A. Halpin, Steven J. Phipps, Joanne M. Whittaker} (2020). \reftitle{The Antarctic Crust and Upper Mantle: A Flexible 3D Model and Software Framework for Interdisciplinary Research}. \textit{Frontiers in Earth Science}, 8, 577502. Frontiers Media SA.\\
\href{https://doi.org/10.3389/feart.2020.577502}{doi:10.3389/feart.2020.577502}
\refabstract{Interdisciplinary research concerning solid Earth–cryosphere interaction and feedbacks requires a working model of the Antarctic crust and upper mantle. Active areas of interest include the effect of the heterogeneous Earth structure on glacial isostatic adjustment, the distribution of geothermal heat, and the history of erosion and deposition. In response to this research need, we construct an adaptable and updatable 3D grid model in a software framework to contain and process solid Earth data. The computational framework, based on an open source software package agrid , allows different data sources to be combined and jointly analyzed. The grid model is populated with crustal properties from geological observations and geochronology results, where such data exist, and published segmentation from geophysical data in the interior where direct observations are absent. The grid also contains 3D geophysical data such as wave speed and derived temperature from seismic tomographic models, and 2D datasets such as gravity anomalies, surface elevation, subglacial temperature, and ice sheet boundaries. We demonstrate the usage of the framework by computing new estimates of subglacial steady-state heat flow in a continental scale model for east Antarctica and a regional scale model for the Wilkes Basin in Victoria Land. We hope that the 3D model and framework will be used widely across the solid Earth and cryosphere research communities.}
\end{refentry}

\begin{refentry}{staal2020grid}
\refauthors{Tobias Stål, Anya M. Reading} (2020). \reftitle{A Grid for Multidimensional and Multivariate Spatial Representation and
                        Data Processing}. \textit{Journal of Open Research Software}, 8(1), 2. Ubiquity Press, Ltd..\\
\href{https://doi.org/10.5334/jors.287}{doi:10.5334/jors.287}
\end{refentry}

\begin{refentry}{staal2019multivariate}
\refauthors{Tobias Stål, Anya M. Reading, Jacqueline A. Halpin, Joanne M. Whittaker} (2019). \reftitle{A Multivariate Approach for Mapping Lithospheric Domain Boundaries in East Antarctica}. \textit{Geophysical Research Letters}, 46(17-18), 10404–10416. American Geophysical Union (AGU).\\
\href{https://doi.org/10.1029/2019GL083453}{doi:10.1029/2019GL083453}
\refabstract{Beneath the ice of East Antarctica lies a continent that is likely to be as geologically complex as its neighbors in Gondwana. An improved model of the heterogeneous lithosphere is required to progress research on Antarctica's tectonic evolution and support interdisciplinary studies of cryosphere and solid Earth interaction. We make use of multiple data sets, which were updated following the field campaigns and compilations of the International Polar Year of 2007/2008. Seismic tomography results, gravity anomalies, and surface elevation are used in a novel method, which combines spatial multivariate data to map possible boundaries as projected likelihood functions. Six multivariate combinations are tested and compared with sparse geological observations in East Antarctica. The resulting lithospheric domain boundaries contribute to our understanding of the deep continental structure. New boundaries are suggested in the interior, and models agree with likely surface expressions of crustal tectonic boundaries exposed along the coast.}
\end{refentry}

\begin{refentry}{morse2019well}
\refauthors{Peter E. Morse, Anya M. Reading, Tobias Stål} (2019). \reftitle{Well-Posed Geoscientific Visualization Through Interactive Color Mapping}. \textit{Frontiers in Earth Science}, 7, 274. Frontiers Media SA.\\
\href{https://doi.org/10.3389/feart.2019.00274}{doi:10.3389/feart.2019.00274}
\end{refentry}

\cvsection{Submitted and in review}

\begin{refentry}{manassero2026mantle}
\refauthors{María Constanza Manassero, Kate Selway, Anya M. Reading, Niam Askey-Doran, J. R. Peacock, F. Ramirez, Tobias Stål} (2026). \reftitle{Magnetotellurics, mantle viscosity, integrated geophysics, glacial isostatic adjustment, mantle water content, East Antarctica}. \textit{Journal of Geophysical Research: Solid Earth}. In review. Working title.
\end{refentry}

\cvsection{Outreach}

\begin{refentry}{staal2024youngminds}
\refauthors{Tobias Stål, Felicity S. McCormack, Anya M. Reading, Niam Askey-Doran, Jacqueline A. Halpin, Mareen Lösing} (2024). \reftitle{Geothermal Heat Shapes the Antarctic Ice Sheet From Below}. \textit{Frontiers for Young Minds}, 11, 1178537. Frontiers Media SA.\\
\href{https://doi.org/10.3389/frym.2023.1178537}{doi:10.3389/frym.2023.1178537}
\refabstract{Antarctica’s ice sheet is constantly on the move, flowing from the deep, frozen interior of the continent toward the ocean, where it melts. In fact, because the oceans are getting warmer, the Antarctic ice sheet melts faster, and therefore, the sea level is rising. However, predicting how the ice sheet will flow differently from place to place is complicated. The landscape beneath the ice sheet helps to control how fast the ice moves. For example, the ice can stick and deform, or slide smoothly across the land under the ice. Naturally occurring heat from inside the Earth can cause the base of the ice sheet to melt and soften so that it flows more easily, sliding on the meltwater formed. The amount of this geothermal heat varies across Antarctica and is difficult to measure. However, scientists with various expertise can collaborate to understand how much heat there is and how it shapes the ice sheet.}
\end{refentry}

\cvsection{Datasets}

\begin{refentry}{fuchs2024global}
\refauthors{Global Heat Flow Data Assessment Group, Sven Fuchs, Florian Neumann, Ben Norden, Elif Balkan-Pazvantoglu, Samah Elbarbary, Alexey Petrunin, Graeme Beardsmore, Robert Harris, Raquel Negrete-Aranda, Jeffrey Poort, Massimo Verdoya, Shaowen Liu, Emma Chambers, Karina Fuentes, Eswara Rao Sidigam, Jhon Camilo Matiz-Leon, Mohammed Hichem Bencharef, Belay G. Mino, Mohamed Shafik Khaled, Denise Verch, Leonard Berger, Saman Firdaus Chishti, Viktoria Dergunova, Helena Liebing, Marvin Schulz, Pia Schuppe, Zlata Trepalova, Paolo Chiozzi, Maria Rosa Alves Duque, Florian Forster, Martina Leveni, Tobias Stål} (2024). \reftitle{The Global Heat Flow Database: Release 2024}. \textit{GFZ Data Services}. Dataset.\\
\href{https://doi.org/10.5880/fidgeo.2024.014}{doi:10.5880/fidgeo.2024.014} \quad \url{https://dataservices.gfz.de/10.5880/fidgeo.2024.014/the-global-heat-flow-database-release-2024}
\refabstract{The data publication contains the compilation of global heat-flow data by the International Heat Flow Commission (IHFC; www.ihfc-iugg.org) of the International Association of Seismology and Physics of the Earth's Interior (IASPEI). The presented data update release 2024 contains data generated between 1939 and 2024 and constitutes the second intermediate update benefiting from the global collaborative assessment and quality control of the Global Heat Flow Database running since May 2021 (http://assessment.ihfc-iugg.org). The data release comprises new original heat-flow data published since April 2023 (the update 2023). It contains 91,182 heat-flow data from 1,586 publications. 57\% of the reported heat-flow values are from the continental domain (n \textasciitilde{} 54,553), while the remaining 43\% are located in the oceanic domain (n \textasciitilde{} 36,692).}
\end{refentry}

\begin{refentry}{reading2024dml}
\refauthors{Anya M. Reading, Tobias Stål, Ian Kelly} (2024). \reftitle{Ice Sheets and Interactions, Dronning Maud Land and Coats Land}. \textit{International Federation of Digital Seismograph Networks}. Dataset.\\
\href{https://doi.org/10.7914/J9CQ-XD60}{doi:10.7914/J9CQ-XD60} \quad \url{https://www.fdsn.org/networks/detail/8G_2024/}
\refabstract{Broadband and node-type sensors deployed with the aim of understanding ice sheets, the ice-bedrock interface zone, and Earth-ice interactions in Dronning Maud Land and Coats Land, Antarctica.}
\end{refentry}

\begin{refentry}{fuchs2023global}
\refauthors{Global Heat Flow Data Assessment Group, Sven Fuchs, Florian Neumann, Ben Norden, Graeme Beardsmore, Paolo Chiozzi, William Colgan, Ana Paulina Anguiano Dominguez, Maria Rosa Alves Duque, Orlando Miguel Ojeda Espinoza, Florian Forster, Andrea Förster, Robert Fröhder, Karina Fuentes, Marek Hajto, Robert Harris, Argo Jõeleht, Helena Liebing, Shaowen Liu, Gwendolin Lüdtke, Mazlan Madon, Raquel Negrete-Aranda, Jeffrey Poort, Itay J. Reznik, Michael Riedel, Frédérique Rolandone, Tobias Stål, Massimo Verdoya, Jyun-Nai Wu} (2023). \reftitle{The Global Heat Flow Database: Update 2023}. \textit{GFZ Data Services}. Dataset.\\
\href{https://doi.org/10.5880/fidgeo.2023.008}{doi:10.5880/fidgeo.2023.008} \quad \url{https://dataservices.gfz.de/10.5880/fidgeo.2023.008/the-global-heat-flow-database-update-2023}
\refabstract{The data publication contains the compilation of global heat-flow data by the International Heat Flow Commission (IHFC; www.ihfc-iugg.org) of the International Association of Seismology and Physics of the Earth's Interior (IASPEI). The presented data update 2023 contains data generated between 1939 and 2022 and constitutes the first intermediate update benefiting from the global collaborative assessment and quality control of the Global Heat Flow Database running since May 2021 (http://assessment.ihfc-iugg.org).}
\end{refentry}

\begin{refentry}{staal2023ice}
\refauthors{Tobias Stål, Anya M. Reading, Shyla Kupis} (2023). \reftitle{Ice Sheets and Interactions, East Antarctica}. \textit{International Federation of Digital Seismograph Networks}. Dataset.\\
\href{https://doi.org/10.7914/74EN-YT90}{doi:10.7914/74EN-YT90} \quad \url{https://www.fdsn.org/networks/detail/5Q_2023/}
\refabstract{Broadband and node-type sensors deployed with the aim of understanding ice sheets, the ice-bedrock interface zone, and Earth-ice interactions in East Antarctica.}
\end{refentry}

\begin{refentry}{staal2021casey}
\refauthors{Tobias Stål, Anya M. Reading, Shyla Kupis, James Newlands} (2021). \reftitle{Casey-Wilkins Adaptable Array}. \textit{International Federation of Digital Seismograph Networks}. Dataset.\\
\href{https://doi.org/10.7914/SN/Z9_2021}{doi:10.7914/SN/Z9\_2021} \quad \url{https://www.fdsn.org/networks/detail/Z9_2021/}
\refabstract{Rapid Deployment Seismic Recorders for Interdisciplinary Antarctic Research. Targeting ice bedrock interface characterisation, seismic wavefield investigations, and method development for ice sheet seismology.}
\end{refentry}

\begin{refentry}{staal2020pangaea}
\refauthors{Tobias Stål, Anya M Reading, Jacqueline A Halpin, Steven J Phipps, Joanne Whittaker} (2020). \reftitle{The Antarctic crust and upper mantle: a flexible 3D model and framework for interdisciplinary research}. \textit{PANGAEA}. Dataset.\\
\href{https://doi.org/10.1594/PANGAEA.918549}{doi:10.1594/PANGAEA.918549} \quad \url{https://doi.pangaea.de/10.1594/PANGAEA.918549}
\refabstract{Example outputs from agrid model.ant\_gia\_dem\_0.tiff is a simplified GIA model, using seismic data to constrain lithosphere thickness and crustal thickness.Aq0.tiff is a steady-state heat flow model with a simple segmented crust, but not including neotectonism. The values in particularly West Antarctica is therefore underestimated.}
\end{refentry}

\begin{refentry}{reading2015cad}
\refauthors{Anya M. Reading, Tobias Stål} (2015). \reftitle{CAD Seismic Deployment, Casey-Davis Region, East Antarctica}. \textit{International Federation of Digital Seismograph Networks}. University of Tasmania. Dataset.\\
\href{https://doi.org/10.7914/NRW9-9825}{doi:10.7914/NRW9-9825} \quad \url{https://www.fdsn.org/networks/detail/9V_2015/}
\refabstract{A sparse seismic instrument deployment between Casey and Davis stations in East Antarctica, also extending to the east of Casey. Co-located with GNSS stations and magnetotelluric recording sites. The network aims to collect data to inform the Earth response to ice sheet change.}
\end{refentry}

\cvsection{Software and code}

\begin{refentry}{staal2022ripley}
\refauthors{Tobias Stål} (2022). \reftitle{Fast Calculation of Ripley's K and L Functions on a Sphere}. \textit{Zenodo}. Software.\\
\href{https://doi.org/10.5281/zenodo.6865005}{doi:10.5281/zenodo.6865005} \quad \url{https://github.com/TobbeTripitaka/ripley}
\refabstract{Fast code to calculate Ripley's K and L functions for a sphere (e.g. a planet)}
\end{refentry}

\begin{refentry}{staal2019agrid}
\refauthors{Tobias Stål} (2019). \reftitle{agrid: a grid for multidimensional and multivariate spatial modelling}. \textit{Zenodo}. Software.\\
\href{https://doi.org/10.5281/zenodo.2553966}{doi:10.5281/zenodo.2553966} \quad \url{https://github.com/TobbeTripitaka/agrid}
\refabstract{This is a first realise of agrid to start testing and for further development. Expect further 0.x releases during the coming weeks after initial tests.}
\end{refentry}

\cvsection{Conference abstracts}

\begin{refentry}{staal2026aq2}
\refauthors{Tobias Stål, Felicity S. McCormack, Anya M. Reading, Magued Al-Aghbary, Aq2 Interdisciplinary Research Collaboration} (2026). \reftitle{Aq2 and Kq2, refined geothermal heat flow models from multivariate observables for application to ice sheet modelling}. \textit{EGU General Assembly 2026}, EGU26-6138. Copernicus GmbH.\\
\href{https://doi.org/10.5194/egusphere-egu26-6138}{doi:10.5194/egusphere-egu26-6138}
\refabstract{Geothermal heat plays an important role in basal ice sheet processes, potentially influencing the stability of Antarctica’s interior ice sheets. Modelled geothermal heat flow can yield discrepancies due to methodological choices and limited data. Recent advances in multivariate analysis and empirical methods have resolved many of these inconsistencies, yielding more consistent outputs that now serve as valuable inputs to ice sheet simulations. Nevertheless, substantial uncertainties remain, particularly in regions with sparse observational coverage. A key factor in addressing these challenges is the spatial resolution of heat flow maps, which governs how basal melt is represented in ice sheet dynamics. We introduce two new geothermal heat flow models designed for glaciated regions: Aq2, developed for continental Antarctica, and Kq2, developed for Greenland. Both models utilise a common framework and employ a multivariate, empirical similarity approach that integrates 18 of the most recent and highest quality observables with the latest reference geothermal heat database, to which we apply weighting and pre-processing to improve representation. Compared to previous empirical models, Aq2 and Kq2 offer reduced uncertainty, greater robustness, and refined spatial resolution. Geothermal heat flow is mapped onto a 0.5 × 0.5 km grid using a forward redistribution approach, which enables higher spatial resolution by leveraging refined observations where available. The models are openly shared in interoperable formats, complete with uncertainty estimates and reproducible code.}
\end{refentry}

\begin{refentry}{staal2025seismicity}
\refauthors{Tobias Stål, Anya M. Reading, Niam Askey-Doran, Sue Cook, Jakob Gradl, Thomas Hudson, Ian Kelly, Bernd Kulessa, Shyla Kupis, Mareen Lösing, Jared Magyar, Maria Constanza Manassero, Matthias Scheiter, Kate Selway, Sarah Thompson, Ross Turner} (2025). \reftitle{Seismicity of Denman Glacier: Constraints on Geometry and Dynamics}. \textit{EGU General Assembly 2025}, EGU25-5279. Copernicus GmbH.\\
\href{https://doi.org/10.5194/egusphere-egu25-5279}{doi:10.5194/egusphere-egu25-5279}
\refabstract{Denman Glacier is one of the largest outlet glaciers in Antarctica. Despite its potential significance for sea level change, its geometry and dynamics remain poorly constrained, making predictions of its response to the changing climate challenging. Seismic signals arise from internal stress and interaction with the subglacial landscape, however, seismic methods are still an often underutilised tool for investigating glacial characteristics. In fact, seismology provides one of the few tools for directly observing the processes and properties that control outlet glacier stability and flow. We present an overview of the seismic data acquired from a transect across Denman Glacier during the Denman Terrestrial Campaign (2023/24). The transect is located 50 km upstream from the grounding line, where the glacier is 13 km wide, and modelling studies suggest a very deep subglacial trough. We employ a combination of passive and active seismic techniques to study the glacier’s geometry, internal structure, and dynamics. Our findings reveal constraints on the glacier's geometry, providing an independent view of the depth extent of the main trough. The recorded seismicity also offers preliminary insights into the glacier's dynamics, the baseline for the monitoring of changing environmental forcing.}
\end{refentry}

\begin{refentry}{staal2023entropy}
\refauthors{Tobias Stål, Anya M. Reading, Matthew J. Cracknell, Jörg Ebbing, Jacqueline A. Halpin, Ian D. Kelly, Emma J. MacKie, Mohamed Sobh, Ross J. Turner, Joanne M. Whittaker} (2023). \reftitle{Using information entropy to optimise and communicate certainty of continental scale tectonic models}. \textit{EGU General Assembly 2023}, EGU23-4074. Copernicus GmbH.\\
\href{https://doi.org/10.5194/egusphere-egu23-4074}{doi:10.5194/egusphere-egu23-4074}
\refabstract{Antarctic subglacial properties impact geothermal heat, subglacial sedimentation, and glacial isostatic adjustment; critical parameters for predicting the ice sheet's response to warming oceans. However, the tectonic architecture of the Antarctic interior is unresolved, with results dependent on datasets or extrapolation used. Most existing deterministic suggestions are derived from qualitative observations and often presented as robust results; however, they hide possible alternative interpretations. Using information entropy as a measure of certainty, we present a robust tectonic segmentation model generated from similarity analysis of multiple geophysical and geological datasets. The use of information entropy provides us with an unbiased and transparent metric to communicate the ambiguities from the uncertainties of qualitative classifications. Information theory also allows us to test and optimise the methods and data to evaluate how choices impact the distribution of alternative output maps. We further discuss how this metric can quantify the predictive power of parameters as a function of regions with different tectonic settings.}
\end{refentry}

\cvsection{Theses}

\begin{refentry}{staal2021phd}
\refauthors{Tobias Stål} (2021). \reftitle{The Antarctic lithosphere revealed by multivariate analysis}. \textit{University of Tasmania}. PhD thesis.\\
\url{https://figshare.utas.edu.au/articles/thesis/The_Antarctic_lithosphere_revealed_by_multivariate_analysis/23254094}
\end{refentry}

\begin{refentry}{staal2015msc}
\refauthors{Tobias Stål} (2015). \reftitle{High-resolution S- and P-wave seismic mapping of near-surface chalk deposits}. \textit{University of Copenhagen}. MSc thesis.
\end{refentry}

\begin{refentry}{staal2011bsc}
\refauthors{Tobias Stål} (2011). \reftitle{Processes and development of the Faroe North Slide Complex}. \textit{University of Copenhagen}. BSc thesis.
\end{refentry}


\end{document}
